%% -*- coding: utf-8 -*-
\documentclass{beamer}

%% -*- coding: utf-8 -*-
\usetheme{Boadilla} % default
\useoutertheme{infolines}
\setbeamertemplate{navigation symbols}{} 

\usepackage{etex}
\usepackage{alltt}
\usepackage{pifont}
\usepackage{color}
\usepackage[utf8]{inputenc}
%\usepackage{german}
\usepackage{listings}
\definecolor{keywordcolor}{rgb}{0.7,0.1,0.1}
\definecolor{tacticcolor}{rgb}{0.0,0.1,0.6}
\definecolor{commentcolor}{rgb}{0.4,0.4,0.4}
\definecolor{symbolcolor}{rgb}{0.0,0.1,0.6}
\definecolor{sortcolor}{rgb}{0.1,0.5,0.1}
\definecolor{attributecolor}{rgb}{0.7,0.1,0.1}

\def\lstlanguagefiles{lstlean.tex}
\lstset{language=Lean}
\lstset{sensitive=true}
\usepackage{hyperref}
\hypersetup{colorlinks=true}
\usepackage[final]{pdfpages}
\usepackage{url}
\usepackage{arydshln} % dashed lines
\usepackage{tikz}
\usepackage{mathpartir}

\DeclareUnicodeCharacter{3BB}{\ensuremath{\lambda}}


\newcommand\cmark{\ding{51}}
\newcommand\xmark{\ding{55}}

\newcommand{\nat}{\mathbf{N}}

\usepackage[all]{xy}

%% new arrow tip for xy
\newdir{|>}{!/4.5pt/@{|}*:(1,-.2)@^{>}*:(1,+.2)@_{>}}

\newcommand\cid[1]{\textup{\textbf{#1}}} % class names
\newcommand\kw[1]{\textup{\textbf{#1}}}  % key words
\newcommand\tid[1]{\textup{\textsf{#1}}} % type names
\newcommand\vid[1]{\textup{\texttt{#1}}} % value names
\newcommand\Mid[1]{\textup{\texttt{#1}}} % method names

\newcommand\TODO[1][]{{\color{red}{\textbf{TODO: #1}}}}

\newcommand\String[1]{\texttt{\dq{}#1\dq{}}}

\newcommand\ClassHead[1]{%
  \ensuremath{\begin{array}{|l|}
      \hline
      \cid{#1}
      \\\hline
    \end{array}}}
\newcommand\AbstractClass[2]{%
  \ensuremath{\begin{array}{|l|}
      \hline
      \cid{\textit{#1}}
      \\\hline
      #2
      \hline
    \end{array}}}
\newcommand\Class[2]{%
  \ensuremath{\begin{array}{|l|}
      \hline
      \cid{#1}
      \\\hline
      #2
      \hline
    \end{array}}}
\newcommand\Attribute[3][black]{\textcolor{#1}{\Param{#2}{#3}}\\}
\newcommand\Methods{\hline}
\newcommand\MethodSig[3]{\Mid{#2} (#3): \,\tid{#1}\\}
\newcommand\CtorSig[2]{\Mid{#1} (#2)\\}
\newcommand\AbstractMethodSig[3]{\Mid{\textit{#2}} (#3): \,\tid{#1}\\}
\newcommand\Param[2]{\vid{#2}:~\tid{#1}}

\lstset{%
  frame=single,
  xleftmargin=2pt,
  stepnumber=1,
  numbers=left,
  numbersep=5pt,
  numberstyle=\ttfamily\tiny\color[gray]{0.3},
  belowcaptionskip=\bigskipamount,
  captionpos=b,
  escapeinside={*'}{'*},
  % language=java,
  tabsize=2,
  emphstyle={\bf},
  commentstyle=\mdseries\it,
  stringstyle=\mdseries\rmfamily,
  showspaces=false,
  showtabs=false,
  keywordstyle=\bfseries,
  columns=fullflexible,
  basicstyle=\footnotesize\CodeFont,
  showstringspaces=false,
  morecomment=[l]\%,
  rangeprefix=////,
  includerangemarker=false,
}

\newcommand\CodeFont{\sffamily}

\definecolor{lightred}{rgb}{0.8,0,0}
\definecolor{darkgreen}{rgb}{0,0.5,0}
\definecolor{darkblue}{rgb}{0,0,0.5}

\newcommand\highlight[1]{\textcolor{blue}{\emph{#1}}}
\newcommand\GenClass[2]{\cid{#1}\texttt{<}\cid{#2}\texttt{>}}

\newcommand\Colored[3]{\alt<#1>{\textcolor{#2}{#3}}{#3}}

\newcommand\nt[1]{\ensuremath{\langle#1\rangle}}

\newcommand{\free}{\operatorname{free}}
\newcommand{\bound}{\operatorname{bound}}
\newcommand{\var}{\operatorname{var}}
\newcommand\VSPBLS{\vspace{-\baselineskip}}

\newcommand\IF{\textit{IF}}
\newcommand\TRUE{\textit{TRUE}}
\newcommand\FALSE{\textit{FALSE}}

\newcommand\IFZ{\textit{IF0}}
\newcommand\ZERO{\textit{ZERO}}
\newcommand\SUCC{\textit{SUCC}}
\newcommand\ADD{\textit{ADD}}
\newcommand\SUB{\textit{SUB}}
\newcommand\MULT{\textit{MULT}}
\newcommand\DIV{\textit{DIV}}

\newcommand\PAIR{\textit{PAIR}}
\newcommand\FST{\textit{FST}}
\newcommand\SND{\textit{SND}}

\newcommand\CASE{\textit{CASE}}
\newcommand\LEFT{\textit{LEFT}}
\newcommand\RIGHT{\textit{RIGHT}}

\newcommand\Encode[1]{\lceil#1\rceil}
\newcommand\Reduce{\stackrel\ast\rightarrow_\beta}

\newcommand\Nat{\textit{Nat}}
\newcommand\Bool{\textit{Bool}}
\newcommand\Pair{\textit{Pair}}
\newcommand\Tfun[1]{#1\to}

\newcommand\Tenv{A}
\newcommand\Lam[1]{\lambda#1.}
\newcommand\App[1]{#1\,}
\newcommand\Succ{\textit{SUCC}\,}
\newcommand\Let[2]{\textit{let}\,#1=#2\,\textit{in}\,}

\newcommand\calE{\mathcal{E}}
\newcommand\calU{\mathcal{U}}
\newcommand\calP{\mathcal{P}}
\newcommand\calW{\mathcal{W}}

\newcommand\GEN{\textit{gen}}
\newcommand\EFV[1]{\textit{fv} (#1)}
\newcommand\Dom[1]{\textit{dom} (#1)}

%%% Local Variables: 
%%% mode: latex
%%% TeX-master: nil
%%% End: 

%%% frontmatter
%% -*- coding: utf-8 -*-

\title{Functional Programming / Lean}
\subtitle{Introduction}

\author[Peter Thiemann]{Prof. Dr. Peter Thiemann}
\institute[Univ. Freiburg]{Albert-Ludwigs-Universität Freiburg, Germany}
\date{WS 2026/27}


\subtitle{More about lists}
\usepackage{tikz}


\begin{document}

\begin{frame}
  \titlepage
\end{frame}
%----------------------------------------------------------------------
\begin{frame}[fragile]
  \frametitle{Lists recap}
  \begin{block}{Zero or more values}
\begin{verbatim}
[]   [1]   [true, false]    ["a", "bunch", "of", "flowers"]
\end{verbatim}
  \end{block}
  \begin{block}{All elements have the same type}
\begin{verbatim}
[true, false] -- good
[1, "two", false] -- bad, type error
\end{verbatim}
  \end{block}
  \begin{block}{Order matters}
\begin{verbatim}
[1,2,3] != [3,2,1]
\end{verbatim}
  \end{block}
\end{frame}
%----------------------------------------------------------------------
\begin{frame}[fragile]
  \frametitle{List syntax}
\begin{lstlisting}
(1 :: (2 :: (3 :: []))) 
==
1 :: 2 :: 3 :: [] 
==
[1,2,3]
\end{lstlisting}
  Alternatively use \lstinline|List.cons| and \lstinline|List.nil| as list constructors.
  
  Strings are distinct from characters, but conversion is possible.
\begin{lstlisting}
"Hearts".toList == ['H','e','a','r','t','s']
String.ofList ['H','e','a','r','t','s'] == "Hearts"
\end{lstlisting}
\end{frame}
%----------------------------------------------------------------------
\begin{frame}[fragile]
  \frametitle{Defining a list datatype}
    \begin{block}<+->{The values of type [a] are \dots}
    \begin{itemize}
    \item either \texttt{[]}, the empty list
    \item or \texttt{x::xs} where \texttt{x} has type \texttt{a} and
      \texttt{xs} has type \texttt{[a]} \\
      ``\texttt{::}'' is pronounced ``cons''
    \end{itemize}
  \end{block}
  \begin{block}<+->{Corresponding definition}
\begin{lstlisting}
inductive List (α : Type) where
  | nil : List α
  | cons : α → List α → List α
\end{lstlisting}
    \begin{itemize}
    \item New: \lstinline{List} is a parametric datatype with type
      parameter \lstinline{α}
    \item Many functions on lists are also parametric (i.e., \textbf{polymorphic}) 
  \end{itemize}
  \end{block}  
\end{frame}
%----------------------------------------------------------------------
\begin{frame}[fragile]
  \frametitle{Predefined polymorphic functions on lists}
\begin{lstlisting}
List.length   : List α -> Int
List.append   : List α -> List α -> List α
List.flatten  : List (List α) -> List α
List.take     : Nat -> List α -> List α
List.drop     : Nat -> List α -> List α
List.zip      : List α → List β → List (α × β)

List.map      : (α → β) → List α → List β 
List.filter   : (α → Bool) → List α → List α
\end{lstlisting}
\end{frame}
%----------------------------------------------------------------------
\begin{frame}[fragile]
  \frametitle{Std functions on lists}
  \begin{block}<+->{Functions on specific lists}
\begin{lstlisting}
List.and  : List Bool -> Bool
List.or   : List Bool -> Bool

def words (s : String) : List String :=
  ((s.split Char.isWhitespace).toStringList).filter (· != "")

def lines (s : String) : List String :=
  s.lines.toStringList

def unwords (ss : List String) : String :=
  " ".intercalate ss

def unlines (ss : List String) : String :=
  String.join (ss.map (fun s => s ++ "\n"))
\end{lstlisting}
  \end{block}
\end{frame}

\begin{frame}[fragile]
  \frametitle{Std functions on lists}
  \begin{block}<+->{Overloaded functions on lists}
\begin{lstlisting}
#eval ([1, 2, 3, 4] : List Nat).sum      -- 10
#eval ([1, 2, 3, 4] : List Nat).prod     -- 24
#eval ([1, 2, 3, 4] : List Nat).elem 3   -- true
#eval ([3, 1, 4, 2] : List Nat).mergeSort -- [1, 2, 3, 4]
\end{lstlisting}
  \end{block}
  
\end{frame}
%----------------------------------------------------------------------
%----------------------------------------------------------------------
% \begin{frame}[fragile]
%   \frametitle{Some examples \dots}
%   \begin{itemize}
%   \item append, reverse
%   \item sum, product
%   \item take, drop, splitAt
%   \item zip, unzip
%   \item insert, isort, qsort
%   \item QuickCheck: collect, classify
%   \end{itemize}
% \end{frame}
%----------------------------------------------------------------------
\begin{frame}[fragile]
  \frametitle{Quicksort!}
\begin{lstlisting}
def qsort {α : Type} [LT α] [LE α] [DecidableLT α] [DecidableLE α] : List α → List α
  | [] => []
  | x :: xs =>
    let smaller := xs.filter (fun y ↦ decide (y ≤ x))
    let larger := xs.filter (fun y ↦ decide (y > x))
    qsort smaller ++ [x] ++ qsort larger
\end{lstlisting}
  \dots{} is not accepted as is \dots
\end{frame}
% ----------------------------------------------------------------------
\begin{frame}[fragile]
  \frametitle{Quicksort terminates \dots}
  We have to explain Lean that it terminates (details later)
\begin{lstlisting}
termination_by xs => xs.length
decreasing_by
  all_goals
    simp_wf
    apply Nat.lt_succ_of_le
    calc
      _ ≤ xs.attach.length := List.length_filter_le ..
      _ = xs.length := by simp
\end{lstlisting}
\end{frame}
% ----------------------------------------------------------------------
% \begin{frame}[fragile]
%   \frametitle{An unfortunate QuickCheck --- ghci interaction}
%   \begin{block}<+->{Two properties}
% \begin{verbatim}
% prop_take_drop n xs = take n xs ++ drop n xs == xs
% nonprop_take_drop n xs = drop n xs ++ take n xs == xs
% \end{verbatim}
%   \end{block}
%   \begin{block}<+->{Testing  \dots}
% \begin{verbatim}
% *Main> quickCheck prop_take_drop
% +++ OK, passed 100 tests.
% *Main> quickCheck nonprop_take_drop
% +++ OK, passed 100 tests.
% \end{verbatim}
%   \end{block}
%   \begin{block}<+->{Oops! what went wrong?}
% \begin{verbatim}
% prop_take_drop :: Eq a => Int -> [a] -> Bool
% nonprop_take_drop :: Eq a => Int -> [a] -> Bool
% \end{verbatim}
%     \vspace{-\baselineskip}
%     \begin{itemize}
%     \item The properties have polymorphic types, but\dots
%     \item QuickCheck does not work with polymorphic types!
%     \end{itemize}
%   \end{block}
% \end{frame}
%----------------------------------------------------------------------
% \begin{frame}[fragile]
%   \frametitle{Ghci ``helps''}
%   \begin{itemize}
%   \item Instead of indicating the problem, ghci chooses a more specific \textbf{default type} 
%   \item In this case, it plugs the unit type for \texttt{a}
%   \item QuickCheck tests
% \begin{verbatim}
% prop_take_drop :: Eq a => Int -> [()] -> Bool
% nonprop_take_drop :: Eq a => Int -> [()] -> Bool
% \end{verbatim}
%   \item Order does not matter when all elements are the same\dots
%   \end{itemize}
% \end{frame}
%----------------------------------------------------------------------
% \begin{frame}[fragile]
%   \frametitle{Force ghci to be unhelpful}
%   \begin{itemize}
%   \item Use type signatures
%   \item Disable defaulting
% \begin{verbatim}
% *Main> :set -XNoExtendedDefaultRules
% \end{verbatim}
%   \item Restrict types used in defaulting
% \begin{verbatim}
% *Main> default (Integer, Double)
% \end{verbatim}
%   \end{itemize}
% \end{frame}
%----------------------------------------------------------------------

\begin{frame}
  \frametitle{Break Time --- Questions?}
  \begin{center}
    \tikz{\node[scale=15] at (0,0){?};}
  \end{center}
\end{frame}


\end{document}

%%% Local Variables: 
%%% mode: latex
%%% TeX-master: t
%%% End: 
